\chapter{Representative Architectures for Node-, Edge-, and Graph-Level Tasks}
\label{ch:architectures}
\minitoc

The message-passing framework of Chapter~\ref{ch:message-passing} fixes the form
of a GNN layer without fixing its content: a message per neighbor, an
order-insensitive aggregation, and an update, with the three functions left open.
The design space of Section~\ref{sec:design-space} consists of the admissible
choices for these functions. This chapter selects three points in that space, one for each
of the task levels introduced in Chapter~\ref{ch:applications}. Each architecture
is presented as an instance of the general
equation~\eqref{eq:message-passing}, and each is paired with the elements its task
demands beyond the message-passing layer itself: how a per-node prediction is
formed for node-level tasks, how a pair of nodes is scored for edge-level tasks,
and how the per-node representations are reduced to a single vector for
graph-level tasks.

\section{Node-Level Learning}
\label{sec:node-level-arch}

A node-level task assigns a prediction to each node of a single graph.
Chapter~\ref{ch:message-passing} established that stacking
$K$ layers produces, for every node, an embedding $\mathbf{h}_i^{(K)}$ that
summarizes the node together with its $K$-hop context, and that the per-node
output of an equivariant network follows the nodes under relabeling, as a
node-level task requires. Chapter~\ref{ch:message-passing} also established the
transductive setting in which such a task is trained. What remains is to specify
how a prediction is computed from each embedding, and then to fix the three
functions of the layer.

\subsection{From Embeddings to Predictions}
\label{sec:node-task-elements}

In node classification each node carries a label drawn from a fixed set of
classes, and the goal is to predict the labels of the unlabeled nodes. The
embedding $\mathbf{h}_i^{(K)}$ produced by the final message-passing layer is the
representation on which this prediction is based. A prediction is obtained by
applying a \textit{classifier} to each node's embedding: a function shared across nodes 
that produces a prediction from each embedding. A common choice is a linear classifier, 
which maps the embedding to one score per class through a learnable linear transformation 
and then normalizes these scores to a probability distribution,
\begin{equation}
  \hat{\mathbf{y}}_i = \operatorname{softmax}\!\left(
    \mathbf{W}\,\mathbf{h}_i^{(K)}
  \right),
  \label{eq:node-classifier}
\end{equation}
where $\mathbf{W}$ is a weight matrix with one row per class, the bias understood
as in Section~\ref{sec:general-equation}, and $\operatorname{softmax}$ the
function that turns a vector of scores into nonnegative entries summing to one.
The entry of $\hat{\mathbf{y}}_i$ with largest value is the predicted class of
node $i$. Because the classifier is shared across nodes and the embeddings are
produced equivariantly, the prediction attached to a node is unchanged by a
relabeling of the graph, that is, it moves with the node.

\subsection{The Graph Convolutional Network}
\label{sec:gcn}

The \textit{Graph Convolutional Network (GCN)} is the instance of the message-passing
framework obtained by taking the aggregation to be a normalized sum over the closed
neighborhood, each term weighted by the degrees of its two incident
nodes~\cite{kipf2017}. It is presented here first in the per-node form that
exhibits it as an instance of the general equation, and then in the matrix form
that connects it to the spectral perspective of
Section~\ref{sec:laplacian-spectral}.

As an instance of the general message-passing equation, GCN fixes its 
three functions as follows. Messages are aggregated over the 
closed neighborhood $\mathcal{N}(i)\cup\{i\}$ by summing the representations
 of the node and its neighbors, each weighted by a normalization coefficient, so that
the aggregation is a normalized sum. The update then applies a shared linear transformation 
followed by an elementwise nonlinearity. The resulting layer computes
\begin{equation}
  \mathbf{h}_i^{(k)} = \sigma\!\left(
    \mathbf{W}^{(k)}
    \sum_{j \in \mathcal{N}(i) \cup \{i\}}
      \frac{1}{\sqrt{\widetilde{\deg}(i)}\,\sqrt{\widetilde{\deg}(j)}}\,
      \mathbf{h}_j^{(k-1)}
  \right),
  \label{eq:gcn-node}
\end{equation}
where $\mathbf{W}^{(k)}$ is the shared weight matrix,
$\sigma$ is an elementwise nonlinearity, and
$\widetilde{\deg}(v)=\deg(v)+1$ denotes the degree of any node $v$
after adding a self-loop.

Every term of the sum ---
including the node's own, entering through the self-loop with $j=i$ --- is weighted
by the coefficient $1/\!\sqrt{\widetilde{\deg}(i)\widetilde{\deg}(j)}$. This
\textit{symmetric normalization} scales each contribution by the degrees of the two
incident nodes, preventing the aggregated representation from growing with the size
of the neighborhood and reducing the influence of highly connected nodes. 

Because the sum runs over the closed neighborhood rather than over
$\mathcal{N}(i)$ alone, the separate update of the central node used in the general
framework is not performed as a distinct step. The node and its neighbors are
folded into a single normalized aggregation and processed by the same linear
transformation.

Applying the same operation to all nodes yields the matrix form of the layer.
The self-loops and the degree normalization are captured by two matrices.
Adding a self-loop at every node gives
\begin{equation}
  \tilde{\mathbf{A}} = \mathbf{A} + \mathbf{I},
  \qquad
  \tilde{D}_{ii} = \sum_{j} \tilde{A}_{ij},
  \label{eq:gcn-renormalization}
\end{equation}
where $\tilde{\mathbf{A}}$ is the adjacency matrix with self-loops and
$\tilde{\mathbf{D}}$ is its degree matrix, the diagonal matrix holding the row
sums of $\tilde{\mathbf{A}}$. With these matrices the layer is
\begin{equation}
  \mathbf{H}^{(k)} = \sigma\!\left(
    \tilde{\mathbf{D}}^{-1/2}\,
    \tilde{\mathbf{A}}\,
    \tilde{\mathbf{D}}^{-1/2}\,
    \mathbf{H}^{(k-1)}\,
    \mathbf{W}^{(k)}
  \right),
  \label{eq:gcn-matrix}
\end{equation}
where $\mathbf{H}^{(k)}$ is the matrix of node representations of
Section~\ref{sec:message-passing-layer}. The matrix
$\tilde{\mathbf{D}}^{-1/2}\tilde{\mathbf{A}}\tilde{\mathbf{D}}^{-1/2}$ that acts on
$\mathbf{H}^{(k-1)}$ is the \textit{propagation operator} of the layer, the object that
mixes each node's representation with those of its neighbors. Its factors
$\tilde{\mathbf{D}}^{-1/2}$ on each side of $\tilde{\mathbf{A}}$ supply the
symmetric normalization, reproducing the coefficient
$1/\!\sqrt{\widetilde{\deg}(i)}\sqrt{\widetilde{\deg}(j)}$
of the per-node form for each pair.
 
The propagation operator connects the layer to the spectral perspective of
Section~\ref{sec:laplacian-spectral}. In the original formulation, GCN is obtained
by applying a first-order approximation to a spectral graph convolution,
yielding the operator $\mathbf{I} + \mathbf{D}^{-1/2}\mathbf{A}\mathbf{D}^{-1/2}$
acting on the node representations before the learnable linear transformation. In
terms of the symmetric normalized Laplacian of
Equation~\eqref{eq:normalized-laplacian}, this operator equals
$2\mathbf{I} - \mathbf{L}_{\mathrm{sym}}$. This is the sense in which the layer is
\textit{spectral}: it acts through a fixed function of the graph Laplacian, the operator 
whose spectrum Section~\ref{sec:laplacian-spectral} showed to be tied to the connectivity of the
graph. The intermediate steps of this spectral derivation are beyond the present
scope, which uses only the resulting operator; the complete derivation is given in
the original work~\cite[-5pc]{kipf2017}.
 
\textit{The renormalization trick}~\eqref{eq:gcn-renormalization} is what turns this
first-order operator into the propagation operator the layer actually applies. It
addresses a spectral difficulty: the eigenvalues of
$\mathbf{I} + \mathbf{D}^{-1/2}\mathbf{A}\mathbf{D}^{-1/2}$ lie in the interval
$[0, 2]$, so applying it repeatedly, as a deep stack of layers does,
can amplify the node representations to the point of numerical instability~\cite{kipf2017}.
Substituting $\tilde{\mathbf{A}}$ for $\mathbf{A}$ and $\tilde{\mathbf{D}}$ for the
degree matrix of $\mathbf{A}$ shrinks this spectrum, reducing the largest eigenvalue
to well below $2$ in practice, which stabilizes the repeated
application~\cite[1pc]{wu2019}.
 
The spectral reading also explains the effect a GCN layer has on the
representations. The quadratic form $\mathbf{f}^\top \mathbf{L} \mathbf{f}$ of
Equation~\eqref{eq:laplacian-quadratic-form} measures how much an assignment of
values to nodes varies across edges, and the same notion of variation applies to
$\mathbf{L}_{\mathrm{sym}}$: an eigenvector of
$\mathbf{L}_{\mathrm{sym}}$ with a large eigenvalue is one whose values vary most
across edges, and one with a small eigenvalue varies little. The propagation
operator applies a fixed function of this Laplacian that inverts the order of its
spectrum, sending the large eigenvalues to small factors and the small eigenvalues
to large ones, so that the components varying most across edges are attenuated and
those varying least are preserved. It thus acts as a low-pass filter on the
graph~\cite{wu2019}, and repeated application reduces the variation across edges,
bringing the representations of adjacent nodes closer together with every layer.
This is \textit{smoothing}: a GCN layer smooths the node representations over the
edges of the graph.

Smoothing is the operative bias of GCN, and it is the right bias exactly when the
graph is homophilous. Section~\ref{sec:structural-properties} described a graph as
homophilous when adjacent nodes tend to share attributes.
On such a graph, making the representations of adjacent nodes
more similar aligns the learned representations with the labels, which is why GCN is effective
for node classification. The same mechanism 
becomes a liability when the assumption fails --- when adjacent nodes
differ systematically, or when so many layers are stacked that all representations
are smoothed toward a common value --- a limitation taken up in
Chapter~\ref{ch:limitations}.

\section{Edge-Level Learning}
\label{sec:edge-level-arch}

An edge-level task makes predictions about pairs of nodes rather than nodes
individually. The representative task, introduced in
Chapter~\ref{ch:applications}, is link prediction: given the edges observed in a
graph, predict which further edges are likely to exist. A common approach to this 
setting is provided by the encoder--decoder framework, and requires two
elements absent from the node-level case: a way to score a pair of nodes from
their embeddings, and a way to construct a supervised training signal from a graph
whose only data is the presence or absence of edges.

\subsection{The Encoder--Decoder Framework}
\label{sec:encoder-decoder}

The encoder-decoder framework divides a link-prediction model into two stages
(Figure~\ref{fig:encoder-decoder}). The
\textit{encoder} is a message-passing network of the kind built in
Chapter~\ref{ch:message-passing}: it maps the graph to node embeddings,
\begin{equation}
  \mathbf{Z} = \operatorname{ENC}(\mathbf{X}, \mathbf{A}),
  \label{eq:encoder}
\end{equation}
where $\mathbf{Z} \in \mathbb{R}^{n \times d'}$ collects the final embeddings, its
row $i$ being the embedding $\mathbf{z}_i = \mathbf{h}_i^{(K)}$ of node $i$. The
encoder carries all the graph structure into the embeddings; what remains is to
turn a pair of embeddings into a prediction about the edge between them. The
\textit{decoder} performs this second step. It is a function that maps a pair of
embeddings to a scalar score expressing how plausible an edge between the two nodes
is. A suitable choice is the \textit{inner product decoder}~\cite{kipf2016}, 
which scores a pair by the inner product of their embeddings 
passed through a function that maps it to a probability,
\begin{equation}
  \hat{A}_{ij} = \sigma\!\left( \mathbf{z}_i^\top \mathbf{z}_j \right),
  \label{eq:inner-product-decoder}
\end{equation}
where $\sigma$ is the logistic function, mapping the real-valued score to a value
in $(0,1)$ read as the predicted probability of an edge between $i$ and $j$. The
inner product is large when the two embeddings point in a similar direction, so the
decoder assigns high probability to pairs whose embeddings the encoder has placed
close together. The decoder requires no parameters of its own, which places the full
responsibility for the prediction on the embeddings.

\begin{figure*}[t]
  \centering
  \includegraphics[width=\textwidth]{encoder-decoder-link-prediction}
  \caption[The encoder--decoder framework for link prediction.]{%
    The encoder--decoder framework for link prediction. The encoder, a message-passing
    network, maps the graph $(\mathbf{X}, \mathbf{A})$ to node embeddings
    $\mathbf{Z}$, one vector $\mathbf{z}_i$ per node. The decoder takes a pair of
    embeddings $\mathbf{z}_i, \mathbf{z}_j$ and returns a scalar score
    $\hat{A}_{ij} = \sigma(\mathbf{z}_i^\top \mathbf{z}_j)$, the predicted
    probability of an edge between $i$ and $j$. A high score is assigned where the
    encoder has placed the two embeddings close together.
  }
  \label{fig:encoder-decoder}
\end{figure*}

The encoder must therefore arrange the embeddings so that connected nodes are close and 
unconnected ones far apart. To train it toward this arrangement, a supervised signal is needed, 
and the structure of the graph itself supplies it. The edges actually present in the graph are the
\textit{positive examples}: pairs that should receive a high score. Pairs of nodes
with no observed edge between them are the \textit{negative examples}: pairs that should
receive a low score. A link-prediction model is trained to raise the score on
positive pairs and lower it on negative pairs, exactly as a binary classifier is
trained to separate two classes.

Two practical points govern how these examples are formed. First, the positive edges are split, 
as in any supervised problem: a portion is held out for evaluation and removed from $\mathbf{A}$ 
during training, not merely from the loss, since otherwise the encoder would see the very edges it is
meant to predict. Second, the negative examples must be sampled, because non-edges vastly outnumber 
observed edges: a graph has at most $\binom{n}{2}=n(n-1)/2$ possible edges but typically 
far fewer actual ones. The standard remedy is \textit{negative sampling}, which draws a random subset 
of non-edges, comparable in number to the positive edges, so that the training signal remains balanced.

\subsection{The Graph Autoencoder}
\label{sec:gae}

The \textit{Graph Autoencoder (GAE)} is the instance of the encoder--decoder framework that
pairs a GCN encoder with the inner product decoder~\cite{kipf2016}. Its name comes
from the autoencoder pattern: a model that encodes its input into a latent
representation and then reconstructs the input from that representation. Here the
input reconstructed is the adjacency matrix. The encoder is a GCN, producing the
embedding matrix $\mathbf{Z}$ as in~\eqref{eq:encoder}; the decoder applies the
inner product~\eqref{eq:inner-product-decoder} to every pair of embeddings at once,
reconstructing the full adjacency matrix as
\begin{equation}
  \hat{\mathbf{A}} = \sigma\!\left( \mathbf{Z} \mathbf{Z}^\top \right),
  \label{eq:gae-reconstruction}
\end{equation}
where $\sigma$ is applied elementwise and the entry $\hat{A}_{ij}$ is the predicted
probability of an edge between $i$ and $j$.

Training matches the reconstructed adjacency against the observed one on the
positive and sampled negative pairs: the GCN
encoder is fitted so that the embeddings it produces, when combined by the inner
product, reproduce the edges of the graph. The arrangement the encoder learns is
the one the decoder rewards --- embeddings of connected nodes made similar, those
of unconnected nodes made dissimilar --- which is the arrangement a GCN
produces by construction, since its smoothing brings the embeddings of adjacent
nodes together.

\section{Graph-Level Learning}
\label{sec:graph-level-arch}

A graph-level task assigns a single prediction to an entire graph. The
representative task, introduced in Chapter~\ref{ch:applications}, is graph
classification: assigning a label to a whole graph, as in predicting a property of
a molecule. This setting requires an element not needed in the previous levels --- a
way to reduce the per-node embeddings, which vary in number from graph to graph,
to a single fixed-size vector representing the graph --- and it places a demand on the
message-passing layer that the previous tasks did not, since graph classification requires 
the layer to preserve the structural information that distinguishes one graph from another.

\subsection{Pooling and Readout}
\label{sec:pooling-readout}

A message-passing network produces one embedding per node. A graph-level
prediction requires one vector for the whole graph, so the per-node embeddings must
be combined into a single graph embedding. This reduction is the \textit{readout}
operation named in Section~\ref{sec:general-equation}, also called \textit{global
pooling}: a function that maps the collection of node embeddings to one vector,
\begin{equation}
  \mathbf{h}_{\mathcal{G}}
    = \operatorname{READOUT}\!\left(
        \{\, \mathbf{h}_i^{(K)} : i \in V \,\}
      \right),
  \label{eq:readout}
\end{equation}
where $\mathbf{h}_{\mathcal{G}}$ is the graph embedding, on which a classifier of
the form~\eqref{eq:node-classifier} is then applied to predict the graph's label.

The readout is subject to one essential requirement:
a graph-level prediction must be unchanged by a
relabeling of the nodes, so the readout must be invariant to the order of its
inputs. It takes a collection of node embeddings and must return the same vector whatever
order they are listed in. This is the permutation invariance of
Equation~\eqref{eq:permutation-invariance},
and it restricts the readout to the order-insensitive
operations of Section~\ref{sec:aggregation-function}: the sum, the mean, and the
maximum, taken coordinatewise over the node embeddings. The readout 
is neighborhood aggregation applied once more, over all nodes at once.

The choice among sum, mean, and maximum is not neutral. The mean
and the maximum can discard information that distinguishes graphs: the mean of a set is
unchanged when the set is duplicated, and the maximum ignores all but the extreme
value in each coordinate, so two graphs differing in the multiplicity or the
distribution of their node embeddings can be mapped to the same vector. The sum
retains more, because it grows with the number of nodes and reflects the full
collection rather than its average or its extremes. This difference is decisive for
the architecture that follows, whose design begins from the requirement that no
distinguishing information be discarded.

\subsection{The Graph Isomorphism Network}
\label{sec:gin}

The \textit{Graph Isomorphism Network (GIN)} is the instance of the 
message-passing framework designed to retain as much distinguishing information 
as the framework allows~\cite{xu2019}. Its motivation is the expressive-power question
introduced in Section~\ref{sec:expressive-power}: a graph-level model is only as good as 
its ability to tell non-isomorphic graphs apart, and, as established there, no
message-passing network can exceed the discriminative power of the \mbox{1-WL} test. 
GIN is the architecture that reaches this ceiling, and the requirement that fixes its design is 
the preservation of neighborhood information through the aggregation.

The requirement can be stated precisely with one notion not yet needed. A
neighborhood, taken as the collection of its neighbors' representations, may
contain the same representation more than once, since distinct neighbors may carry
identical vectors. The right object is therefore not a set, which records only
which elements are present, but a \textit{multiset}, a collection in which elements
may repeat and the number of repetitions is part of the object. For the aggregation
to preserve all the information in a neighborhood, it must be \textit{injective}
over multisets: it must map distinct neighborhood multisets to distinct vectors,
never collapsing two different neighborhoods to the same aggregate. The mean and
the maximum fail this requirement, for the reasons discussed previously.
The sum, by contrast, can be made injective
over multisets, because it reflects every element and its multiplicity.

GIN builds its layer on the sum accordingly. The aggregation is the sum over the
neighborhood; the update combines the node's own representation with this sum and
passes the result through a function expressive enough to preserve the
distinction, giving
\begin{equation}
  \mathbf{h}_i^{(k)}
    = \operatorname{MLP}^{(k)}\!\left(
        \left( 1 + \epsilon^{(k)} \right) \mathbf{h}_i^{(k-1)}
        + \sum_{j \in \mathcal{N}(i)} \mathbf{h}_j^{(k-1)}
      \right),
  \label{eq:gin}
\end{equation}
where $\operatorname{MLP}^{(k)}$ is the update function of
the layer, and $\epsilon^{(k)}$ is a scalar, which can be fixed or learnable. 
The MLP in the update can approximate an injective function of
its input, so the composition of the sum and the update preserves the distinction
between different closed neighborhoods. The term $(1 + \epsilon^{(k)})$ keeps the
node's own representation distinguishable from the aggregate of its neighbors,
rather than letting the two merge as they do in the closed-neighborhood sum of GCN.
Together these choices make one GIN layer as discriminating as one round of the
\mbox{1-WL} refinement of Section~\ref{sec:expressive-power}, which is the most a
message-passing layer can be.